According to a theorem of \'Elie Cartan, a smooth map $f \colon \Sigma \rightarrow G$, from a
connected manifold $\Sigma$ into a Lie group $G$, is uniquely
determined by its {\it logarithmic derivative}, up to right
translations in $G$. This derivative, also known as the Darboux
derivative of $f$, is a one-form $\delta f$ on~$\Sigma $ taking values
in the Lie algebra of ${\mathfrak g} $ of $G$. Here we formulate and
prove a generalization of this result, Theorem \ref{summaryT}, to
smooth maps $f \colon \Sigma \rightarrow G/H$ into an arbitrary
homogeneous space $G/H$. Our generalization describes explicitly the
global obstruction to reconstructing such maps from infinitesimal
data, data that generalizes logarithmic derivatives (generalized
Maurer--Cartan forms).

In this introduction we generalize the notion of Maurer--Cartan forms
and their monodromy, and state the main existence Theorem~\ref{outlineT}. The proof is straightforward, apart from a question
about Lie algebroid integrability, which is addressed in Section
\ref{newsec}, by applying \cite{Crainic_Fernandes_03}. Uniqueness, up
to symmetry, is guaranteed under a mild topological condition on
$G/H$, but we must take some care to qualify what is meant by
``symmetry'', a task postponed to Section~\ref{infch}.
